\chapter{Resultados e discussão}
\label{chap:resultados}

Este capítulo apresenta os resultados na ordem do \textit{pipeline}: as bacias hidrográficas (\autoref{sec:res-bacias}), a calibração do CCA e as manchas urbanas (\autoref{sec:res-cca}), a repartição das manchas entre as bacias (\autoref{sec:res-reparticao}), o microplástico associado às fozes e o efeito do raio de aceitação (\autoref{sec:res-plastico}), as leis de escala entre população e microplástico (\autoref{sec:res-alometria}) e as verificações das distribuições (\autoref{sec:res-distribuicoes}). A \autoref{sec:discussao} discute o conjunto dos resultados e as suas limitações.

\section{Bacias hidrográficas}
\label{sec:res-bacias}

O relevo reamostrado tem 11.851 \textit{sinks} costeiros, e o IPBA produziu o mesmo número de bacias, uma por \textit{sink}. Ordenadas por área, as 138 maiores cobrem 90,0\% da área de terra da região de estudo (\autoref{fig:bacias}); as outras 11.713 são, na grande maioria, micro-bacias de encosta com poucas células. As bacias principais têm área mediana de 13.572 km$^2$, e a menor delas tem 4.897 km$^2$. Seis bacias principais tocam a borda da caixa ou o corte do Canal do Panamá e estão marcadas como truncadas; entre elas estão as duas maiores, a do Mississippi e a do São Lourenço, cujas cabeceiras mais ao norte ficam acima de 50$^\circ$N.

\begin{figure}[ht!]
	\centering
	\Caption{\label{fig:bacias}As 138 bacias principais, que cobrem 90\% da área de terra da região de estudo}
	\IBGEtab{}{%
		\includegraphics[width=\textwidth]{bacias_principais}%
	}{%
		\Fonte{elaborado pelo autor.}%
		\Nota{Cada bacia principal tem uma cor (as cores se repetem); as bacias costeiras menores estão em cinza-claro.}%
	}
\end{figure}

As bacias delineadas reproduzem as grandes bacias conhecidas do continente sem nenhuma informação além do relevo e da linha de costa. A bacia do Mississippi--Missouri ocupa o centro dos Estados Unidos, do oeste de Montana aos Apalaches; a do São Lourenço inclui os Grandes Lagos; a do Colorado drena o sudoeste até o Golfo da Califórnia; e a do Columbia drena o noroeste. No México aparecem as bacias do Rio Bravo (Rio Grande), do Pánuco, do Lerma--Santiago e do Balsas. A \autoref{tab:bacias-pop} lista as dez bacias principais mais populosas. Os nomes foram atribuídos a partir da coordenada da foz e não entram em nenhum cálculo.

\begin{table}[ht!]
	\centering
	\Caption{\label{tab:bacias-pop}As dez bacias principais mais populosas}
	\IBGEtab{}{%
		\small
		\begin{tabular}{llrrrr}
			\toprule
			\textbf{Bacia} & \textbf{Foz (lon, lat)} & \textbf{Área (km$^2$)} & \textbf{População (mi)} & \textbf{\% urbana} & \textbf{Sítio (km)} \\
			\midrule
			Mississippi--Atchafalaya$^{*}$ & $-91{,}18$; $29{,}58$ & 3.150.557 & 70,20 & 45 & 37,8 \\
			São Lourenço$^{*}$ & $-70{,}88$; $46{,}92$ & 1.374.739 & 61,10 & 60 & 316,9 \\
			Pánuco & $-97{,}85$; $22{,}35$ & 166.272 & 31,50 & 84 & 114,5 \\
			Rio Bravo (Rio Grande) & $-97{,}32$; $26{,}28$ & 849.584 & 24,24 & 84 & 9,7 \\
			Lerma--Santiago & $-105{,}38$; $21{,}58$ & 133.470 & 22,60 & 75 & 92,3 \\
			Balsas & $-102{,}18$; $17{,}95$ & 117.024 & 16,55 & 65 & 260,9 \\
			Colorado & $-114{,}82$; $31{,}85$ & 809.665 & 15,87 & 73 & 229,9 \\
			Columbia$^{*}$ & $-123{,}58$; $46{,}22$ & 1.012.275 & 12,86 & 64 & 35,7 \\
			Sacramento--San Joaquin & $-122{,}25$; $38{,}08$ & 177.542 & 9,85 & 75 & 21,3 \\
			Mobile & $-87{,}98$; $30{,}72$ & 171.579 & 8,62 & 18 & 2,0 \\
			\bottomrule
		\end{tabular}%
	}{%
		\Fonte{elaborado pelo autor.}%
		\Nota{$^{*}$ bacia truncada pela borda da caixa. ``\% urbana'' é a fração da população da bacia que vive em manchas urbanas do CCA. ``Sítio'' é a distância da foz ao sítio do \textit{Nurdle Patrol} mais próximo.}%
	}
\end{table}

A população somada das 138 bacias principais é de 450,4 milhões de habitantes, 74\% dos 606,6 milhões da região de estudo; o restante vive nas bacias costeiras menores, que concentram cidades litorâneas como Nova York e boa parte da Flórida. Na passagem das bacias para a grade do GPW, 8,5 milhões de habitantes de células da linha de costa foram realocados para a bacia de terra mais próxima, e apenas 168.498 habitantes ficaram fora de qualquer bacia.

\section{Calibração do CCA e manchas urbanas}
\label{sec:res-cca}

\subsection{Matriz de expoentes}

A \autoref{fig:zipf-matriz} mostra o expoente $\hat{\alpha}_{100}$ nas 15.000 combinações de $(D^{*}, \ell)$ e o desvio $\Delta$ da \autoref{eq:delta}. Abaixo de cerca de 1.000 m, $\ell$ é menor que o espaçamento entre células na mesma coluna ($\approx 927$ m), e só se ligam células vizinhas na mesma linha em latitudes altas, onde as células são mais estreitas; quase toda célula ativa é uma mancha, e o expoente fica entre 0,3 e 0,6, refletindo sobretudo a distribuição das densidades das células. Entre 1.000 e 1.300 m, as células se ligam na horizontal e na vertical, mas perto do equador ainda não na diagonal, e aparecem as manchas degeneradas descritas na \autoref{sec:calibracao}. A partir de $\ell_{\min} = 1.400$ m, a curva $\hat{\alpha}_{100} = 1$ atravessa a matriz em valores de $D^{*}$ entre cerca de 400 e 1.100 hab/km$^2$.

\begin{figure}[ht!]
	\centering
	\Caption{\label{fig:zipf-matriz}Expoente de Zipf e desvio $\Delta$ na grade de parâmetros do CCA}
	\IBGEtab{}{%
		\includegraphics[width=\textwidth]{zipf_matriz}%
	}{%
		\Fonte{elaborado pelo autor.}%
		\Nota{(a) $\hat{\alpha}_{100}$ em cada combinação de $D^{*}$ e $\ell$, com as curvas $\hat{\alpha}_N = 1$ para $N = 50$, 100 e 200. (b) Desvio $\Delta = \max_N |\hat{\alpha}_N - 1|$. A faixa hachurada, $\ell < \ell_{\min} = 1.400$ m, é excluída da escolha; a estrela marca o par escolhido.}%
	}
\end{figure}

A \autoref{tab:calibracao} mostra, para cada $\ell$ entre 1.400 e 2.600 m, o $D^{*}$ que minimiza $\Delta$. O mínimo global ocorre em $\ell = 1.700$ m e $D^{*} = 1.010$ hab/km$^2$, com $\Delta = 0{,}081$. O mínimo é raso: todos os valores de $\ell$ entre 1.400 e 2.400 m têm um $D^{*}$ com $\Delta < 0{,}10$, e o $D^{*}$ ótimo fica entre 820 e 1.110 hab/km$^2$ em quase todos eles. Em todas as linhas da tabela, $\hat{\alpha}_{200}$ é maior que 1 e é ele que determina $\Delta$; os expoentes do topo, $\hat{\alpha}_{50}$ e $\hat{\alpha}_{100}$, ficam entre 0,90 e 0,97. Esse padrão, em que o expoente cresce com $N$, é discutido na \autoref{sec:res-rank-size}.

\begin{table}[ht!]
	\centering
	\Caption{\label{tab:calibracao}Melhor $D^{*}$ para cada $\ell$ admitido pelo critério $\Delta$}
	\IBGEtab{}{%
		\small
		\begin{tabular}{rrrrrr}
			\toprule
			\textbf{$\ell$ (m)} & \textbf{$D^{*}$ (hab/km$^2$)} & \textbf{$\Delta$} & \textbf{$\hat{\alpha}_{50}$} & \textbf{$\hat{\alpha}_{100}$} & \textbf{$\hat{\alpha}_{200}$} \\
			\midrule
			1.400 & 960 & 0,081 & 0,919 & 0,926 & 1,053 \\
			1.500 & 970 & 0,088 & 0,917 & 0,912 & 1,079 \\
			1.600 & 1.040 & 0,098 & 0,903 & 0,917 & 1,098 \\
			\textbf{1.700} & \textbf{1.010} & \textbf{0,081} & \textbf{0,925} & \textbf{0,939} & \textbf{1,081} \\
			1.800 & 630 & 0,090 & 0,910 & 0,915 & 1,088 \\
			1.900 & 820 & 0,091 & 0,959 & 0,909 & 1,089 \\
			2.000 & 850 & 0,097 & 0,942 & 0,903 & 1,096 \\
			2.100 & 1.110 & 0,090 & 0,936 & 0,945 & 1,090 \\
			2.200 & 1.100 & 0,093 & 0,963 & 0,950 & 1,093 \\
			2.300 & 1.110 & 0,083 & 0,965 & 0,942 & 1,083 \\
			2.400 & 1.110 & 0,085 & 0,952 & 0,947 & 1,085 \\
			2.500 & 1.110 & 0,100 & 0,962 & 0,959 & 1,100 \\
			2.600 & 1.110 & 0,108 & 0,961 & 0,965 & 1,108 \\
			\bottomrule
		\end{tabular}%
	}{%
		\Fonte{elaborado pelo autor.}%
		\Nota{Em negrito, o par escolhido. Para $\ell > 2.600$ m, o menor $\Delta$ passa de 0,13.}%
	}
\end{table}

O par escolhido tem uma interpretação física razoável. O limiar de 1.010 hab/km$^2$ é cerca de 2,6 vezes o limiar de área urbana do censo americano (386 hab/km$^2$) \cite{census2020}, o que é esperado, já que o GPW distribui a população de cada unidade censitária de forma homogênea e suaviza os picos de densidade. A distância de 1,7 km é pouco maior que a diagonal entre células vizinhas, o que significa que as manchas são praticamente conjuntos de células ativas contíguas, com tolerância para vãos de uma célula.

\subsection{Manchas urbanas}

Com $D^{*} = 1.010$ hab/km$^2$ e $\ell = 1.700$ m, o CCA identifica 12.185 manchas urbanas, que somam 349,5 milhões de habitantes, 58\% da população da região. O \texttt{cca\_grid} aplicado com esses parâmetros produziu exatamente o mesmo número de manchas e a mesma população da maior mancha que a varredura (20.602.596 habitantes), confirmando a equivalência entre os dois procedimentos. A \autoref{tab:top-manchas} lista as dez maiores manchas. Sem nenhuma lista de cidades fornecida ao algoritmo, elas correspondem às grandes metrópoles da região, com os centroides a poucos quilômetros dos centros urbanos.

\begin{table}[ht!]
	\centering
	\Caption{\label{tab:top-manchas}As dez maiores manchas urbanas ($D^{*} = 1.010$ hab/km$^2$, $\ell = 1.700$ m)}
	\IBGEtab{}{%
		\small
		\begin{tabular}{rlrrrr}
			\toprule
			\textbf{Posto} & \textbf{Mancha} & \textbf{População} & \textbf{Área (km$^2$)} & \textbf{Densidade (hab/km$^2$)} & \textbf{Centroide (lon, lat)} \\
			\midrule
			1 & Cidade do México & 20.602.596 & 2.204 & 9.349 & $-99{,}10$; $19{,}45$ \\
			2 & Nova York & 15.482.384 & 3.316 & 4.669 & $-73{,}92$; $40{,}75$ \\
			3 & Los Angeles & 14.240.127 & 4.017 & 3.545 & $-118{,}03$; $33{,}98$ \\
			4 & Chicago & 6.335.092 & 2.253 & 2.812 & $-87{,}83$; $41{,}85$ \\
			5 & Toronto & 6.029.628 & 1.372 & 4.394 & $-79{,}49$; $43{,}71$ \\
			6 & Miami & 5.329.254 & 2.073 & 2.570 & $-80{,}24$; $26{,}05$ \\
			7 & Guadalajara & 5.073.091 & 594 & 8.547 & $-103{,}35$; $20{,}63$ \\
			8 & Monterrey & 4.971.276 & 637 & 7.800 & $-100{,}28$; $25{,}72$ \\
			9 & Houston & 4.814.175 & 1.982 & 2.429 & $-95{,}46$; $29{,}76$ \\
			10 & Santo Domingo & 4.737.554 & 947 & 5.004 & $-69{,}92$; $18{,}50$ \\
			\bottomrule
		\end{tabular}%
	}{%
		\Fonte{elaborado pelo autor a partir de \citeonline{gpw4}.}%
	}
\end{table}

A tabela tem duas particularidades. A primeira é que as manchas mexicanas e caribenhas são mais densas que as americanas: Cidade do México, Guadalajara e Monterrey têm mais de 7.800 hab/km$^2$, contra 2.400 a 4.700 hab/km$^2$ nas metrópoles dos Estados Unidos, o que reflete a forma mais compacta dessas cidades. A segunda é um artefato do GPW em Monterrey: cerca de 3 milhões de pessoas estão concentradas em uns 20 pixels, com densidades de até 810 mil hab/km$^2$, fisicamente impossíveis. O artefato provavelmente vem de uma unidade censitária muito pequena que recebeu uma população grande. Ele não muda a mancha, porque todas essas células já estão acima de $D^{*}$ e a população total é a mesma, mas mostra que valores de densidade de células individuais do GPW precisam ser lidos com cautela. O dado não foi alterado.

\subsection{Distribuição \textit{rank-size}}
\label{sec:res-rank-size}

A \autoref{fig:zipf} mostra a distribuição \textit{rank-size} das manchas. No topo do \textit{ranking}, a distribuição segue de perto a lei de Zipf: os estimadores de Hill dão $\hat{\alpha}_{50} = 0{,}925$, $\hat{\alpha}_{100} = 0{,}939$ e $\hat{\alpha}_{200} = 1{,}081$. Mais abaixo no \textit{ranking}, a curva fica mais íngreme. O ajuste CSN sobre a cauda escolhe $x_{\min} = 21.121$ habitantes, com 1.325 manchas na cauda, e dá $\hat{\gamma} = 1{,}822 \pm 0{,}023$, ou seja, $\hat{\alpha} = 1{,}217 \pm 0{,}033$. A distância KS é 0,017, abaixo do valor crítico de 0,037, e o \textit{bootstrap} dá $p = 0{,}08$, próximo do limiar convencional de 0,1. A regressão log-log sobre a mesma cauda dá $\alpha = 1{,}16$ ($R^2 = 0{,}994$).

\begin{figure}[ht!]
	\centering
	\Caption{\label{fig:zipf}Distribuição \textit{rank-size} das 12.185 manchas urbanas}
	\IBGEtab{}{%
		\includegraphics[width=0.85\textwidth]{zipf}%
	}{%
		\Fonte{elaborado pelo autor.}%
		\Nota{Retas de Hill para as 50, 100 e 200 maiores manchas e reta do ajuste CSN sobre a cauda ($x \geq x_{\min}$).}%
	}
\end{figure}

A distribuição é, portanto, \textbf{côncava} no gráfico log-log: o expoente cresce à medida que se incluem manchas menores. A comparação pela razão de Vuong confirma essa leitura. A lei de potência é fortemente preferida à exponencial ($R = 8{,}64$, $p < 10^{-17}$), não se distingue da lognormal ($R = -1{,}29$, $p = 0{,}20$) e perde para a lei de potência truncada ($R = -2{,}52$, $p = 0{,}004$), que tem exatamente essa curvatura. O resultado concorda com a literatura: \citeonline{eeckhout2004} mostrou que a distribuição completa das cidades americanas é lognormal e que a lei de Zipf vale apenas na cauda superior, e \citeonline{rozenfeld2011} encontraram a lei de Zipf nas manchas do CCA acima de um tamanho mínimo. A conclusão prática é que a calibração garante a lei de Zipf para as maiores manchas, que são as que mais pesam na população das bacias, mas não para a distribuição inteira.

\section{Repartição das manchas entre as bacias}
\label{sec:res-reparticao}

Das 12.185 manchas, 618 cruzam divisores de água e foram divididas entre duas ou mais bacias. São poucas em número, mas grandes: somam 41\% da população urbana, o que mostra que a repartição não é um detalhe. A \autoref{fig:manchas-fozes} mostra as dez maiores manchas e as fozes das bacias principais que recebem pedaços delas, e a \autoref{fig:mapa-integrado} sobrepõe as bacias, as manchas e os sítios de coleta.

\begin{figure}[ht!]
	\centering
	\Caption{\label{fig:manchas-fozes}As dez maiores manchas urbanas e as fozes das bacias principais em que se dividem}
	\IBGEtab{}{%
		\includegraphics[width=\textwidth]{mapa_manchas_fozes}%
	}{%
		\Fonte{elaborado pelo autor.}%
		\Nota{Cada mancha tem uma cor e um símbolo; o símbolo marca a foz de cada bacia principal que recebe um pedaço da mancha, ligada ao centroide por uma linha pontilhada. Na exibição, as dez manchas foram dilatadas em cerca de 11 km para ficarem visíveis na escala continental. As demais manchas estão em cinza.}%
	}
\end{figure}

Os casos das maiores manchas ilustram os três comportamentos possíveis.

\begin{itemize}
	\item \textbf{Mancha inteira numa bacia principal.} A Cidade do México está 100\% na bacia do Pánuco, que deságua no Golfo do México perto de Tampico, a mais de 300 km da cidade; o mesmo vale para Guadalajara (Lerma--Santiago), Monterrey (Rio Bravo) e Toronto (São Lourenço). São casos em que o método atribui a população a uma foz distante, que uma análise baseada em proximidade costeira não encontraria.
	\item \textbf{Mancha dividida entre bacias principais.} Chicago fica sobre o divisor entre o Lago Michigan e o rio Illinois: 83\% da mancha drena para o São Lourenço e 17\% para o Mississippi. Houston se divide entre cinco bacias, duas delas principais.
	\item \textbf{Mancha em bacias costeiras pequenas.} Nova York se divide entre 119 bacias, nenhuma delas principal; o maior pedaço tem 10\% da mancha. Com um \textit{sink} por pixel de costa, a Baía de Nova York e a foz do Hudson se fragmentam em muitas micro-bacias na grade de 2$'$, e a cidade não é associada a nenhuma bacia de rio. Los Angeles e Miami têm comportamento semelhante: cerca de 40\% e 56\% das manchas, respectivamente, ficam em bacias principais, e o restante se espalha por dezenas de bacias costeiras.
\end{itemize}

\begin{figure}[ht!]
	\centering
	\Caption{\label{fig:mapa-integrado}Bacias principais, manchas urbanas, sítios de coleta e fozes}
	\IBGEtab{}{%
		\includegraphics[width=\textwidth]{mapa_integrado}%
	}{%
		\Fonte{elaborado pelo autor.}%
		\Nota{Triângulos cheios: fozes com sítio do \textit{Nurdle Patrol} a até $R_{\mathrm{maioria}} = 38$ km; triângulos vazios: fozes sem sítio nesse raio.}%
	}
\end{figure}

O último caso é uma limitação importante. Grandes cidades litorâneas que ficam em estuários fragmentados, como Nova York, são justamente as que mais poderiam contribuir com plástico para o mar e ficam fora das bacias principais. Uma solução seria agrupar os \textit{sinks} de uma mesma baía ou estuário antes de rodar o IPBA, o que é discutido como trabalho futuro no \autoref{chap:conclusao}.

\section{Microplástico nas fozes e raio de aceitação}
\label{sec:res-plastico}

Dentro da caixa ampliada pela margem de 50 km há 5.783 amostras do \textit{Nurdle Patrol} em 4.814 sítios distintos, dos quais 861 têm contagem média zero. A distribuição espacial é muito desigual (\autoref{fig:pontos-coleta}): a maior parte dos sítios está na costa do Golfo do México, do Texas à Flórida, com grupos menores na costa atlântica, na costa do Pacífico e nos Grandes Lagos, e poucos sítios no México, na América Central e no Caribe. A distância mediana entre a foz de uma bacia principal e o sítio mais próximo é de 88,1 km, e apenas 6 das 138 fozes têm um sítio a até 1 km.

A \autoref{fig:varredura} e a \autoref{tab:varredura} mostram o efeito do raio de aceitação. A fração dos sítios a até $R$ de alguma foz principal cresce devagar até cerca de 30 km e acelera depois disso; o menor raio que contém a maioria dos sítios é $R_{\mathrm{maioria}} = 38$ km, com 53\% deles. O número de bacias com sítio e contagem positiva cresce rápido nos primeiros 10 km, chegando a 21, e depois devagar, até 33 em 50 km.

\begin{figure}[ht!]
	\centering
	\Caption{\label{fig:varredura}Efeito do raio de aceitação entre a foz e o sítio de coleta}
	\IBGEtab{}{%
		\includegraphics[width=\textwidth]{varredura_raio}%
	}{%
		\Fonte{elaborado pelo autor.}%
		\Nota{(a) Fração dos sítios a até $R$ de uma foz principal, número de bacias com sítio e número de sítios distintos. (b) Coeficiente de Pearson $r$ entre os logaritmos da população e do microplástico. (c) Expoente $\beta$ com IC de 95\%. Linha tracejada: $R = 1$ km; linha pontilhada: $R_{\mathrm{maioria}} = 38$ km. Só entram as bacias com contagem positiva.}%
	}
\end{figure}

\begin{table}[ht!]
	\centering
	\Caption{\label{tab:varredura}Leis de escala entre população e microplástico em função do raio de aceitação}
	\IBGEtab{}{%
		\small
		\begin{tabular}{rrrrrrrrr}
			\toprule
			& & & & \multicolumn{3}{c}{\textbf{Concentração}} & \multicolumn{2}{c}{\textbf{Conc. $\times$ área}} \\
			\cmidrule(lr){5-7} \cmidrule(lr){8-9}
			\textbf{$R$ (km)} & \textbf{\% sítios} & \textbf{Bacias} & \textbf{Sítios} & \textbf{$\beta$} & \textbf{$r$} & \textbf{$p$} & \textbf{$\beta$} & \textbf{$r$} \\
			\midrule
			1 & 3,3 & 3 & 3 & -- & -- & -- & -- & -- \\
			2 & 5,2 & 6 & 6 & 0,386 & 0,936 & 0,006 & 0,313 & 0,486 \\
			5 & 7,5 & 17 & 17 & 0,443 & 0,427 & 0,088 & 0,672 & 0,511 \\
			10 & 12,1 & 21 & 21 & 0,464 & 0,430 & 0,052 & 0,912 & 0,578 \\
			20 & 24,8 & 25 & 24 & 0,306 & 0,363 & 0,075 & 0,683 & 0,562 \\
			30 & 37,1 & 27 & 26 & 0,318 & 0,376 & 0,054 & 0,699 & 0,567 \\
			\textbf{38} & \textbf{53,0} & \textbf{29} & \textbf{28} & \textbf{0,283} & \textbf{0,367} & \textbf{0,050} & \textbf{0,849} & \textbf{0,665} \\
			40 & 61,0 & 30 & 29 & 0,267 & 0,355 & 0,054 & 0,816 & 0,655 \\
			50 & 82,8 & 33 & 31 & 0,243 & 0,339 & 0,054 & 0,778 & 0,652 \\
			\bottomrule
		\end{tabular}%
	}{%
		\Fonte{elaborado pelo autor.}%
		\Nota{``Bacias'': bacias principais com sítio a até $R$ e contagem positiva; ``Sítios'': sítios distintos entre elas. $p$ é o valor-$p$ do coeficiente de Pearson. Com $R = 1$ km há apenas 3 bacias com contagem positiva, insuficientes para o ajuste.}%
	}
\end{table}

\textbf{Medição garantida ($R = 1$ km).} Seis bacias principais têm um sítio a até 1 km da foz, mas três delas têm contagem zero, e restam apenas três com contagem positiva: as de Los Angeles (San Gabriel), a de Galveston--Houston e uma pequena bacia do sul do Texas. Com três pontos não é possível ajustar uma lei de escala nem avaliar a sua significância. O painel (a) da \autoref{fig:dispersao} mostra esses pontos apenas para referência.

\textbf{Raios pequenos.} Com $R = 2$ km, as seis bacias com contagem positiva dão $r = 0{,}94$ ($p = 0{,}006$) para a concentração. O valor é alto, mas se apoia em seis pontos e cai para cerca de 0,43 já em $R = 5$ km, com 17 bacias. Ele é um exemplo do risco de escolher o raio pelo resultado: a varredura mostra que valores altos de $r$ em raios pequenos são flutuações de amostras minúsculas.

\textbf{Estabilidade a partir de 10 km.} De 10 a 50 km, o $r$ da concentração fica entre 0,34 e 0,45 e o $\beta$ entre 0,24 e 0,47, com valores-$p$ entre 0,05 e 0,09. Para o indicador de quantidade, $r$ fica entre 0,56 e 0,67 e $\beta$ entre 0,68 e 0,91. A relação é, portanto, estável em relação ao raio, mas fraca e no limite da significância estatística convencional para a concentração.

\section{Leis de escala entre população e microplástico}
\label{sec:res-alometria}

A \autoref{fig:dispersao} mostra as relações entre a população das bacias e o microplástico nas fozes, e a \autoref{tab:escala} resume os ajustes para $R_{\mathrm{maioria}} = 38$ km. As 29 bacias com contagem positiva usam 28 sítios distintos.

\begin{figure}[ht!]
	\centering
	\Caption{\label{fig:dispersao}População das bacias versus microplástico na foz}
	\IBGEtab{}{%
		\includegraphics[width=\textwidth]{dispersao_pop_lixo}%
	}{%
		\Fonte{elaborado pelo autor.}%
		\Nota{(a) e (b): bacias com sítio a até 1 km da foz. (c) e (d): bacias com sítio a até $R_{\mathrm{maioria}} = 38$ km. Em (a) e (c), concentração de \textit{pellets}; em (b) e (d), concentração multiplicada pela área da bacia. Retas: MQO em log-log; faixas: banda de confiança de 95\% por \textit{bootstrap}.}%
	}
\end{figure}

\begin{table}[ht!]
	\centering
	\Caption{\label{tab:escala}Leis de escala entre população da bacia e microplástico na foz ($R = 38$ km)}
	\IBGEtab{}{%
		\small
		\begin{tabular}{lrrrrrrr}
			\toprule
			\textbf{Variável} & \textbf{$n$} & \textbf{$\beta$} & \textbf{IC 95\%} & \textbf{$p_{\mathrm{perm}}$} & \textbf{$r$} & \textbf{$\rho_s$} & \textbf{$R^2$} \\
			\midrule
			Concentração & 29 & 0,283 & [0,108; 0,539] & 0,049 & 0,367 & 0,480 & 0,135 \\
			Concentração $\times$ área & 29 & 0,849 & [0,503; 1,321] & $<0{,}001$ & 0,665 & 0,684 & 0,442 \\
			\bottomrule
		\end{tabular}%
	}{%
		\Fonte{elaborado pelo autor.}%
		\Nota{IC por \textit{bootstrap} de pares (10.000 reamostragens); $p_{\mathrm{perm}}$ por permutação (10.000 permutações).}%
	}
\end{table}

\textbf{Concentração.} A concentração de \textit{pellets} na foz cresce com a população da bacia segundo um expoente $\beta = 0{,}28$, com IC de 95\% de 0,11 a 0,54, que exclui o zero, e $p_{\mathrm{perm}} = 0{,}049$. A relação é fortemente \textbf{sublinear}: dobrar a população da bacia está associado a um aumento de apenas $2^{0{,}28} - 1 \approx 21\%$ na concentração. A dispersão é grande: $r = 0{,}37$ e $R^2 = 0{,}13$, ou seja, a população explica cerca de 13\% da variância do logaritmo da concentração. A correlação de postos é um pouco maior, $\rho_s = 0{,}48$ ($p = 0{,}008$), o que indica que parte da fraqueza do ajuste linear vem de pontos extremos, como a bacia com 471 \textit{pellets} por 10 minutos na costa central do Texas.

\textbf{Indicador de quantidade.} Multiplicando a concentração pela área da bacia, o expoente sobe para $\beta = 0{,}85$ (IC de 0,50 a 1,32), com $r = 0{,}66$ e $p_{\mathrm{perm}} < 0{,}001$. O resultado parece muito mais forte, mas a maior parte dele vem da área, e não do plástico. Como $\log(\text{concentração} \times A) = \log(\text{concentração}) + \log A$, e o MQO é linear, o expoente do indicador de quantidade é a soma do expoente da concentração com o expoente da área em função da população: $0{,}85 = 0{,}28 + 0{,}57$. Nas bacias com sítio, a área e a população estão fortemente correlacionadas ($\rho_s = 0{,}67$, $p < 10^{-7}$), porque bacias maiores tendem a ter mais gente. O indicador de quantidade herda essa correlação, que existiria mesmo que a concentração de plástico fosse aleatória. Ele não deve, portanto, ser lido como evidência independente de que a população determina o plástico.

\textbf{Contagens zero.} Das 52 bacias principais com sítio a até 38 km, 23 têm contagem zero e ficam fora dos ajustes em log-log. A exclusão não é neutra: se as bacias com zero fossem sistematicamente menos populosas, excluí-las aumentaria artificialmente a inclinação. Os dados não indicam esse viés: a população mediana das bacias com zero é de 1,10 milhão de habitantes, contra 1,97 milhão nas bacias com contagem positiva, diferença que não é significativa (teste de Mann--Whitney, $p = 0{,}50$). Mas a inclusão dos zeros enfraquece a associação. Com as 52 bacias, a correlação de postos entre população e concentração cai para $\rho_s = 0{,}25$ ($p = 0{,}07$), e entre população e indicador de quantidade para $\rho_s = 0{,}32$ ($p = 0{,}02$).

\textbf{Influência de bacias individuais.} O teste \textit{leave-one-out} (\autoref{fig:loo}) mostra que nenhuma bacia isolada determina o resultado da concentração: removendo uma bacia de cada vez, $r$ varia de 0,33 a 0,43 e $\beta$ de 0,25 a 0,35. As bacias de maior influência são as mais populosas, Mississippi e Rio Bravo, cuja remoção reduz $r$ em cerca de 0,035, e uma bacia do noroeste da Flórida, com foz perto de Panama City, cuja remoção aumenta $r$ em 0,06.

\begin{figure}[ht!]
	\centering
	\Caption{\label{fig:loo}Variação do coeficiente $r$ ao remover cada bacia do ajuste da concentração ($R = 38$ km)}
	\IBGEtab{}{%
		\includegraphics[width=0.6\textwidth]{leave_one_out}%
	}{%
		\Fonte{elaborado pelo autor.}%
		\Nota{A faixa cinza marca $|\Delta r| \leq 0{,}1$.}%
	}
\end{figure}

\section{Distribuições de tamanho}
\label{sec:res-distribuicoes}

A \autoref{fig:leis-potencia} e a \autoref{tab:leis-potencia} mostram os ajustes de cauda das áreas das bacias, das populações das bacias e das populações das manchas urbanas. As três distribuições têm cauda pesada, e em todas a lei de potência é fortemente preferida à exponencial.

\begin{figure}[ht!]
	\centering
	\Caption{\label{fig:leis-potencia}Distribuições complementares acumuladas das áreas e populações das bacias e das populações das manchas urbanas}
	\IBGEtab{}{%
		\includegraphics[width=\textwidth]{leis_potencia}%
	}{%
		\Fonte{elaborado pelo autor.}%
		\Nota{Linhas: lei de potência (contínua) e lognormal (tracejada) ajustadas acima de $x_{\min}$ (linha vertical pontilhada). $R_{\mathrm{LN}}$ é a razão de verossimilhança normalizada entre a lei de potência e a lognormal.}%
	}
\end{figure}

\begin{table}[ht!]
	\centering
	\Caption{\label{tab:leis-potencia}Ajustes de cauda pelo método CSN}
	\IBGEtab{}{%
		\small
		\begin{tabular}{lrrrrrr}
			\toprule
			\textbf{Distribuição} & \textbf{$n$} & \textbf{$x_{\min}$} & \textbf{Cauda} & \textbf{$\hat{\gamma}$} & \textbf{$R_{\mathrm{LN}}$ ($p$)} & \textbf{$R_{\mathrm{exp}}$ ($p$)} \\
			\midrule
			Área das bacias (km$^2$) & 11.794 & 82 & 1.877 & $1{,}655 \pm 0{,}015$ & $-1{,}20$ (0,23) & 8,19 ($<10^{-15}$) \\
			População das bacias & 11.134 & 86.142 & 482 & $1{,}795 \pm 0{,}036$ & $-0{,}54$ (0,59) & 5,93 ($<10^{-8}$) \\
			População das manchas & 12.185 & 7.301 & 2.939 & $1{,}782 \pm 0{,}014$ & $-1{,}76$ (0,08) & 10,99 ($<10^{-27}$) \\
			\bottomrule
		\end{tabular}%
	}{%
		\Fonte{elaborado pelo autor.}%
		\Nota{Área das bacias: apenas as que não tocam a borda. $R_{\mathrm{LN}}$ e $R_{\mathrm{exp}}$: razões de verossimilhança normalizadas entre a lei de potência e a lognormal e a exponencial; valores positivos favorecem a lei de potência.}%
	}
\end{table}

A área das bacias segue uma lei de potência com $\hat{\gamma} = 1{,}66$ a partir de 82 km$^2$, indistinguível da lognormal. O expoente é maior que os valores de 1,4 a 1,5 típicos da distribuição de áreas de sub-bacias aninhadas numa rede fluvial \cite{rodrigueziturbe1997,maritan1996}, o que é esperado: aqui as bacias não são aninhadas, mas uma partição do continente em bacias costeiras disjuntas, e as bacias grandes ``roubam'' área das vizinhas. A população das bacias tem cauda com $\hat{\gamma} = 1{,}80$, também indistinguível da lognormal. A população das manchas urbanas, com $x_{\min}$ mais baixo que no ajuste do primeiro \textit{notebook}, dá $\hat{\gamma} = 1{,}78$, ou seja, $\alpha = 1/(\hat{\gamma} - 1) \approx 1{,}28$, coerente com a concavidade discutida na \autoref{sec:res-rank-size}; nesse caso, a lognormal é marginalmente preferida ($p = 0{,}08$).

\section{Discussão}
\label{sec:discussao}

\subsection{O que os resultados dizem sobre a hipótese}

A hipótese de trabalho era que a quantidade de microplástico na foz cresce com a população da bacia. Os resultados a sustentam apenas em parte. A associação entre população e concentração de \textit{pellets} é positiva em todos os raios testados, estável entre 10 e 50 km e com IC que exclui o zero em $R_{\mathrm{maioria}}$. Mas é fraca, com a população explicando cerca de 13\% da variância, fica no limite da significância estatística e enfraquece quando as contagens zero são incluídas. O indicador de quantidade dá uma relação mais forte, mas herda a correlação entre população e área e não constitui evidência independente.

O expoente sublinear da concentração, $\beta \approx 0{,}28$, tem pelo menos quatro explicações possíveis, que não se excluem.

\begin{enumerate}
	\item \textbf{Diluição.} Bacias mais populosas tendem a ser maiores e a ter vazões maiores. Se a carga de plástico cresce com a população, mas a vazão cresce com a área, a concentração na foz cresce menos que a carga. A concentração numa praia não é uma carga, e uma relação linear entre população e carga é compatível com uma relação sublinear entre população e concentração.
	\item \textbf{Fonte industrial.} Os \textit{pellets} são um microplástico primário cuja perda está ligada à produção e ao transporte de resinas, e não ao consumo \cite{karlsson2018}. As maiores contagens estão na costa do Texas \cite{tunnell2020}, onde se concentra a indústria petroquímica, e não necessariamente nas bacias mais populosas. A população da bacia é, nesse caso, um indicador indireto da atividade industrial.
	\item \textbf{Transporte costeiro.} Uma vez no mar, os \textit{pellets} são redistribuídos por correntes e pela deriva litorânea e depositados em praias que podem estar a dezenas de quilômetros da foz de origem. O sítio mais próximo pode refletir a soma de várias fozes.
	\item \textbf{Ruído amostral.} A maioria dos sítios associados às bacias tem uma única amostra, feita num único dia, e as contagens de \textit{pellets} variam muito com o tempo, as marés e as tempestades.
\end{enumerate}

\subsection{O papel das escolhas metodológicas}

Algumas escolhas feitas ao longo do trabalho protegem os resultados de interpretações otimistas. Os parâmetros do CCA foram fixados pela lei de Zipf antes de qualquer contato com os dados de plástico, e a varredura mostrou que o critério escolhe uma região estável da matriz de parâmetros e não um ponto isolado. O raio de aceitação não foi escolhido pelo resultado, mas por um critério de cobertura dos sítios, e a varredura completa mostra que o resultado não depende criticamente dele. O $r = 0{,}94$ obtido com $R = 2$ km, que seria o resultado mais ``favorável'', foi reportado como o que é: uma flutuação de seis pontos. As contagens zero, que precisam sair do ajuste em log-log, foram analisadas à parte com uma medida de postos.

Outras escolhas limitam o alcance dos resultados:

\begin{itemize}
	\item \textbf{Resolução do relevo.} A reamostragem para 2$'$, imposta pela implementação em Python puro do IPBA, faz com que a foz seja conhecida com precisão de cerca de 3,7 km e fragmenta estuários complexos, como o de Nova York, em muitas micro-bacias. Um ``sítio a até 1 km'' significa a até 1 km desse pixel.
	\item \textbf{Bacias truncadas.} A caixa corta as cabeceiras mais ao norte das bacias do Mississippi e do São Lourenço e as bacias que drenam para a Baía de Hudson; o IPBA atribui as partes dentro da caixa às bacias vizinhas.
	\item \textbf{Cobertura das coletas.} O \textit{Nurdle Patrol} concentra-se no Golfo do México. Das 138 bacias principais, apenas 52 têm sítio a até 38 km, e o México, a América Central e o Caribe contribuem com poucas delas. As bacias com medição não são uma amostra aleatória das bacias do continente.
	\item \textbf{Tipo de microplástico.} Os \textit{pellets} são um tipo específico de microplástico primário, e os resultados não se estendem automaticamente aos microplásticos secundários, como fibras e fragmentos, que têm fontes difusas e mais diretamente ligadas ao consumo.
	\item \textbf{Concentração versus carga.} A contagem por 10 minutos de busca numa praia é uma medida de concentração local, não da carga transportada pelo rio. Um teste mais direto da hipótese exigiria medições de carga fluvial, que não estão disponíveis numa escala continental.
\end{itemize}
